\section{Threat Report Sets the Bound}
\sectionkicker{SAFETY BOUND}
\begin{wideflow}
{\Large\bfseries 安全の境界を、事例の範囲で}\par
\smallskip
{\color{SurveyMuted}9月10日の脅威報告を、今週の技術的な安全の目安として、目立つ事例の留保ごと置く。}\par
\end{wideflow}

9月10日、Anthropicは8か月分（2025年12月から2026年8月）の脅威報告を示した\cite{threat}。扱うharmは7領域で、対象の模型はHaiku、Sonnet、Opusに限られ、蒸留の1件を除きFableやMythos級は関わらないとされる。複雑な攻撃が必ずしも熟練を要さず、複数代理人の枠組みで自律の度が増し、人が狙いを定めて抜き取りを確かめる構図が示される。報告書の日付自体は通常の範囲の出来事として扱う。

取り上げられるのは、ロシアの諜報作戦、場当たりの窃取集団、脆弱性の探索工房、AI供給連鎖の窃取と転売などで、いずれも妨害と守りの強化、当局への共有が添えられる\cite{threat}。生物の研究に関わる5件も妨害した利用の中に数えられるが、兵器が生まれたとか意図が固まったという主張はしていない。いずれも提供側の調査報告であり、帰属は公開報道と整合的とされる範囲で受け止め、ここで独立に確かめない。全文のPDFや検出子の村は今回たどっておらず、抄録の範囲でしか読まない。9月11日の時系列にある蒸留の名指しの細目は範囲後の速報として切り分け、通常の事実には入れない。目立つ事例は全体の縮図ではないという留保を保つ。

\begin{claimboundary}[安全の読み方の境界]
全文や検出子はたどっていない。帰属は提供側の調査報告である。速報の細目は通常の事実に入れない。
\end{claimboundary}
